\section{Conclusion}
Prime Agent introduces a new paradigm for agent harness design in which persistent execution, recursive sessions, autonomous controls, recorded history, and Continual Harness form one substrate for long-horizon work. Results across interactive reasoning, long-context tasks, autonomous research, systems construction, and persistent environments show that this substrate supports different forms of test-time computation under standardized execution and accounting. Despite its results relative to alternative harnesses, models still experience friction when deciding how to allocate subagents, manage retained information, and refine reusable state. Many harness capabilities remain underused because current models were not trained to operate them. We expect model-harness co-learning to become the dominant route to new long-horizon capabilities. Training directly with Prime Agent could teach models to use the integrated harness more effectively, while targeted training on the RLM and Continual Harness components could isolate their contributions.
